\documentclass[12pt,a4paper]{article}

% ─── Packages ────────────────────────────────────────────────────────────────
\usepackage[utf8]{inputenc}
\usepackage[T1]{fontenc}
\usepackage[spanish,english]{babel}
\usepackage{geometry}
\geometry{top=2.5cm, bottom=2.5cm, left=3cm, right=3cm}
\usepackage{amsmath, amssymb, amsthm}
\usepackage{booktabs}
\usepackage{array}
\usepackage{longtable}
\usepackage{fancyhdr}
\usepackage{setspace}
\usepackage{graphicx}
\usepackage{hyperref}
\usepackage{enumitem}
\usepackage{parskip}
\usepackage{caption}
\usepackage{multirow}
\usepackage{xcolor}
\usepackage{titlesec}
\usepackage{tocloft}
\usepackage{natbib}

% ─── Hyperref ────────────────────────────────────────────────────────────────
\hypersetup{
    colorlinks=true,
    linkcolor=black,
    citecolor=black,
    urlcolor=blue,
    pdftitle={SARIMA and SARIMAX: Dataset Selection and Methodological Reference},
    pdfauthor={Guillermo},
    pdfsubject={Time Series Econometrics}
}

% ─── Page Style ──────────────────────────────────────────────────────────────
\pagestyle{fancy}
\fancyhf{}
\rhead{\small Series de Tiempo — Finanzas Cuantitativas}
\lhead{\small Universidad Panamericana}
\cfoot{\thepage}
\renewcommand{\headrulewidth}{0.4pt}

% ─── Section Formatting ──────────────────────────────────────────────────────
\titleformat{\section}[block]{\large\bfseries\scshape}{}{0em}{}[\titlerule]
\titleformat{\subsection}[block]{\normalsize\bfseries}{}{0em}{}
\titleformat{\subsubsection}[block]{\normalsize\itshape}{}{0em}{}

% ─── Theorem Environments ────────────────────────────────────────────────────
\newtheorem{definition}{Definition}[section]
\newtheorem{remark}{Remark}[section]

% ─── Spacing ─────────────────────────────────────────────────────────────────
\setstretch{1.3}

% ═════════════════════════════════════════════════════════════════════════════
\begin{document}
% ═════════════════════════════════════════════════════════════════════════════

% ─── Title Page ──────────────────────────────────────────────────────────────
\begin{titlepage}
\centering
\vspace*{2cm}
{\scshape\Large Universidad Panamericana \par}
{\scshape\large Facultad de Ciencias Empresariales \par}
{\scshape\large Finanzas Cuantitativas — Series de Tiempo \par}
\vspace{2cm}
{\Huge\bfseries SARIMA and SARIMAX:\\[0.4em]
\large Notation, Methodology, Dataset Selection,\\ and Academic Reference \par}
\vspace{2cm}
\rule{0.5\textwidth}{0.4pt}\\[0.5cm]
{\large Guillermo \par}
{\large 8.º Semestre — Área Cuantitativa \par}
\vspace{1cm}
{\large Curso: Dr.\ César Galindo \par}
\vspace{2cm}
{\large \today \par}
\vfill
\end{titlepage}

% ─── Table of Contents ───────────────────────────────────────────────────────
\tableofcontents
\newpage

% ═════════════════════════════════════════════════════════════════════════════
\section{Introduction}
% ═════════════════════════════════════════════════════════════════════════════

The family of Seasonal Autoregressive Integrated Moving Average models — commonly
written as (S)ARIMA$(p,d,q)(P,D,Q)_s$ — constitutes one of the central
frameworks in the econometric analysis of time series with periodic structure.
Introduced in systematic form by \citet{BoxJenkins1976}, extended by
\citet{BoxJenkinsReinsel2015}, and enriched with exogenous regressors in the
SARIMAX specification, these models serve both descriptive and predictive
purposes across macroeconomics, finance, climatology, and the natural sciences.

This document has three purposes. First, it decodes the notation that appears on
the classroom whiteboard — every symbol, every subscript, every diagnostic
criterion — and identifies the most defensible interpretation of the parsimony
constant~$C$ in the constraint $p+d+q+P+D+Q\le C$. Second, it surveys fifteen
publicly available datasets ranked by their fitness for class-level SARIMA and
SARIMAX analysis. Third, it provides a methodological critique of how attached
files and classroom images should be presented in future prompt engineering
sessions, specific to this exact type of econometric assignment.

% ═════════════════════════════════════════════════════════════════════════════
\section{Decoding the Classroom Notation}
% ═════════════════════════════════════════════════════════════════════════════

\subsection{The General Specification}

The whiteboard displays the notation:
\[
    (S)\text{ARIMA}(p,d,q)(P,D,Q)_s
\]
with the following annotations:
\begin{itemize}[noitemsep]
  \item ACF (Autocorrelation Function) — pointing to $q$
  \item PACF (Partial Autocorrelation Function, \textit{esperanza condicional}) — pointing to $p$
  \item ``diff'' — pointing to $d$
  \item ``Seasonality'' — pointing to the entire prefix $(S)$
  \item ``Si aplicamos acf al AR no da información''
  \item On the right board: $p+d+q+P+D+Q\le C$ — \textit{Principio de parsimonia}
  \item Tests: AIC, BIC, log-verosimilitud, Ljung-Box, SSE
\end{itemize}

\subsection{Component Definitions}

\begin{definition}[Non-seasonal order $p$]
The autoregressive order. The current observation is regressed on its own $p$
previous values. Identified via the PACF: significant spikes at lags $1,\dots,p$
and cutoff thereafter indicate an AR$(p)$ structure.
\end{definition}

\begin{definition}[Integration order $d$]
The number of non-seasonal differences required to render the series weakly
stationary. Determined through visual inspection of the ACF (slow decay signals
non-stationarity) and confirmed by unit-root tests such as Augmented
Dickey--Fuller (ADF) or KPSS. A value of $d=1$ is sufficient for most
macroeconomic series; $d=2$ is rare and must be justified.
\end{definition}

\begin{definition}[Non-seasonal MA order $q$]
The moving-average order. The current innovation is modelled as a weighted sum
of $q$ past innovations (white-noise shocks). Identified via the ACF of the
differenced series: significant spikes at lags $1,\dots,q$ with abrupt cutoff
indicate MA$(q)$.
\end{definition}

\begin{definition}[Seasonal AR order $P$]
Autoregressive structure operating at the seasonal lag $s$. If $s=12$ (monthly),
the seasonal AR term at order $P=1$ links $y_t$ to $y_{t-12}$. Identified from
spikes in the PACF at multiples of $s$.
\end{definition}

\begin{definition}[Seasonal integration order $D$]
Number of seasonal differences applied: $\Delta_s y_t = y_t - y_{t-s}$. Removes
deterministic or stochastic seasonality. Typically $D\in\{0,1\}$.
\end{definition}

\begin{definition}[Seasonal MA order $Q$]
Moving-average structure at seasonal lags. Identified from spikes in the ACF of
the seasonally differenced series at multiples of $s$.
\end{definition}

\begin{definition}[Seasonal period $s$]
The fundamental frequency of recurrence: $s=12$ for monthly data with annual
cycles, $s=4$ for quarterly, $s=7$ for daily data with weekly cycles, $s=52$
for weekly data, etc. Must be specified a priori based on domain knowledge.
\end{definition}

\subsection{The Complete Model Equation}

Let $B$ denote the backshift operator ($By_t = y_{t-1}$, $B^s y_t = y_{t-s}$).
Define:
\[
    \phi_p(B) = 1 - \phi_1 B - \phi_2 B^2 - \cdots - \phi_p B^p
\]
\[
    \Phi_P(B^s) = 1 - \Phi_1 B^s - \Phi_2 B^{2s} - \cdots - \Phi_P B^{Ps}
\]
\[
    \theta_q(B) = 1 + \theta_1 B + \theta_2 B^2 + \cdots + \theta_q B^q
\]
\[
    \Theta_Q(B^s) = 1 + \Theta_1 B^s + \Theta_2 B^{2s} + \cdots + \Theta_Q B^{Qs}
\]

The SARIMA$(p,d,q)(P,D,Q)_s$ model is:
\[
    \Phi_P(B^s)\,\phi_p(B)\,(1-B)^d\,(1-B^s)^D\,y_t
    = \Theta_Q(B^s)\,\theta_q(B)\,\varepsilon_t
\]
where $\varepsilon_t\sim\text{WN}(0,\sigma^2)$.

\subsection{The Note on ACF Applied to AR}

The whiteboard annotation \textit{``Si aplicamos acf al AR no da información''}
is a standard teaching point. In a pure AR$(p)$ process, the ACF decays
exponentially (or in a damped sinusoidal pattern) and never cuts off sharply;
it is therefore uninformative for reading $p$ directly. The PACF, which
conditions on all intervening lags — hence \textit{esperanza condicional} —
produces the clean cutoff after lag $p$ that allows direct identification.
Symmetrically, in a pure MA$(q)$ process the ACF cuts off at lag $q$ while the
PACF decays gradually. This is the fundamental Box--Jenkins identification
logic.

% ═════════════════════════════════════════════════════════════════════════════
\section{The Parsimony Constant $C$}
% ═════════════════════════════════════════════════════════════════════════════

\subsection{The Constraint}

The whiteboard states:
\[
    p + d + q + P + D + Q \le C
\]
labelled \textit{Principio de parsimonia}. This is a hard constraint on the
total number of model parameters — or equivalently, on the aggregate order
complexity — used during the grid search over candidate specifications.

\subsection{Most Defensible Interpretation of $C$}

The literature and standard practice converge on the following interpretation:

\begin{center}
\fbox{\parbox{0.85\textwidth}{
\textbf{$C$ is an instructor-specified upper bound on aggregate order complexity,
typically set between 6 and 10 for academic exercises.} Its purpose is to prevent
students from fitting overparameterised models that exploit sample idiosyncrasies
rather than genuine structure.
}}
\end{center}

Three specific benchmark values appear in textbooks and software defaults:
\begin{enumerate}
  \item \textbf{$C = 6$}: Common in introductory courses. Allows, for example,
  SARIMA$(2,1,1)(1,1,1)_{12}$ where the sum is $2+1+1+1+1+1=7>6$, so in
  practice the constraint forces one component to zero. This is intentionally
  restrictive.
  \item \textbf{$C = 8$}: The implicit bound used by Rob Hyndman's \texttt{auto.arima}
  function in R (default \texttt{max.order = 5} per group, effectively capping
  complexity). An $C=8$ version allows slightly richer seasonal models.
  \item \textbf{$C = 10$}: Used in some graduate-level courses where the student
  is expected to explore a wider grid but select parsimoniously via AIC/BIC.
\end{enumerate}

\begin{remark}
In some course handouts, $C$ is defined not as a sum of orders but as the
maximum number of \emph{freely estimated} parameters (i.e., $p + q + P + Q$,
excluding the integration orders $d$ and $D$ which are not estimated but fixed
by differencing). If the professor's constraint includes $d$ and $D$, as the
whiteboard clearly shows, the interpretation is total order complexity rather
than parameter count. For a typical academic dataset with $d=1, D=1$, a constraint
of $C=6$ leaves at most $p+q+P+Q \le 4$, which is a reasonable parsimonious grid.
\end{remark}

% ═════════════════════════════════════════════════════════════════════════════
\section{ARIMA, SARIMA, and SARIMAX: Conceptual Taxonomy}
% ═════════════════════════════════════════════════════════════════════════════

\subsection{ARIMA$(p,d,q)$}

The baseline non-seasonal model. Captures autocorrelation and trends
via differencing. Suitable for series without periodic structure or where
seasonality has been removed by prior adjustment. The process is:
\[
    \phi_p(B)(1-B)^d y_t = \mu + \theta_q(B)\varepsilon_t
\]

\subsection{SARIMA$(p,d,q)(P,D,Q)_s$}

Extends ARIMA with a multiplicative seasonal component. The seasonal
operators $\Phi_P(B^s)$ and $\Theta_Q(B^s)$ model autocorrelation at
multiples of the seasonal period. ``Multiplicative'' refers to the polynomial
multiplication $\Phi_P(B^s)\cdot\phi_p(B)$, which produces cross-lag terms
between seasonal and non-seasonal dynamics. For example, SARIMA$(1,1,1)(1,1,1)_{12}$
— the ``airline model'' of \citet{BoxJenkins1976} — is:
\[
    (1-\Phi B^{12})(1-\phi B)(1-B)(1-B^{12})y_t = (1+\Theta B^{12})(1+\theta B)\varepsilon_t
\]

\subsection{SARIMAX$(p,d,q)(P,D,Q)_s$}

Adds a vector of exogenous regressors $\mathbf{x}_t$:
\[
    \Phi_P(B^s)\phi_p(B)(1-B)^d(1-B^s)^D y_t
    = \boldsymbol{\beta}'\mathbf{x}_t + \Theta_Q(B^s)\theta_q(B)\varepsilon_t
\]
The exogenous variables enter the mean equation rather than the error structure.
They explain variation that the autoregressive and moving-average terms cannot
capture. Examples: trading volume as exogenous to a price series; temperature as
exogenous to electricity demand; oil price as exogenous to an inflation model.
Critically, SARIMAX is not a causal model by construction; exogeneity must be
argued on theoretical grounds, not merely assumed.

% ═════════════════════════════════════════════════════════════════════════════
\section{Diagnostic and Selection Tools}
% ═════════════════════════════════════════════════════════════════════════════

\subsection{Autocorrelation Function (ACF)}

The ACF at lag $k$ is:
\[
    \rho(k) = \frac{\text{Cov}(y_t, y_{t-k})}{\text{Var}(y_t)}
\]
It measures linear dependence between $y_t$ and its own past, ignoring
intermediate lags. In a stationary series, $\rho(k)\to 0$ as $k\to\infty$.
Non-stationary series exhibit ACF decay so slow it appears virtually flat.
The ACF identifies the MA order $q$: it cuts off at lag $q$ for pure MA processes.

\subsection{Partial Autocorrelation Function (PACF)}

The PACF at lag $k$ is the correlation between $y_t$ and $y_{t-k}$ after
removing the linear effect of all intermediate lags $y_{t-1},\dots,y_{t-k+1}$.
This is precisely the \textit{esperanza condicional} (conditional expectation)
referenced on the whiteboard — it corresponds to the coefficient $\hat{\phi}_{kk}$
in the regression of $y_t$ on $(y_{t-1},\dots,y_{t-k})$. The PACF identifies
the AR order $p$: it cuts off at lag $p$ for pure AR processes.

\subsection{Differencing}

First-order differencing: $\Delta y_t = y_t - y_{t-1} = (1-B)y_t$ removes
stochastic trends. Seasonal differencing: $\Delta_s y_t = y_t - y_{t-s} = (1-B^s)y_t$
removes periodic nonstationarity. Over-differencing introduces unnecessary
MA components and should be avoided. The ADF test (and variants) guides the
choice of $d$; visual inspection of the ACF guides $D$.

\subsection{Parsimonious Model Selection}

Parsimony (Occam's razor applied to time series) is not an aesthetic preference
but an epistemological constraint: given finite data, a model with fewer
parameters generalises better out-of-sample than one with more. The constraint
$p+d+q+P+D+Q\le C$ operationalises this by restricting the search grid.
Within that grid, AIC and BIC select among candidates.

\subsection{AIC — Akaike Information Criterion}

\[
    \text{AIC} = -2\,\hat{\ell} + 2k
\]
where $\hat{\ell}$ is the maximised log-likelihood and $k$ is the number of
freely estimated parameters. AIC penalises complexity lightly; it is consistent
for prediction rather than for model identification. Preferred when forecasting
accuracy is the objective.

\subsection{BIC — Bayesian Information Criterion}

\[
    \text{BIC} = -2\,\hat{\ell} + k\ln(n)
\]
The penalty grows with sample size $n$, making BIC more conservative than AIC
for large datasets. BIC is consistent for model selection (it selects the true
model with probability approaching 1 as $n\to\infty$, if the true model is in
the candidate set). Preferred when parsimony is the primary objective.

\subsection{Log-Likelihood (\textit{log-verosimilitud})}

Under Gaussian innovations:
\[
    \hat{\ell}(\boldsymbol{\theta}) = -\frac{n}{2}\ln(2\pi) - \frac{n}{2}\ln(\hat{\sigma}^2)
    - \frac{1}{2\hat{\sigma}^2}\sum_{t=1}^n \hat{\varepsilon}_t^2
\]
Maximising $\hat{\ell}$ over the parameter vector $\boldsymbol{\theta}$ is
equivalent to minimising the sum of squared residuals (SSE) under Gaussian
errors, which explains why SSE also appears on the whiteboard as a selection
criterion. Higher log-likelihood (less negative) indicates better fit.

\subsection{Ljung--Box Test}

After fitting a candidate model, the residuals $\hat{\varepsilon}_t$ should
approximate white noise. The Ljung--Box statistic at lag $m$ is:
\[
    Q_{LB}(m) = n(n+2)\sum_{k=1}^m \frac{\hat{\rho}_k^2}{n-k}
    \;\sim\;\chi^2(m - p - q - P - Q)
\]
A failure to reject ($p$-value $> 0.05$) indicates the residuals are
serially uncorrelated — the model has extracted all exploitable structure.
Rejection implies residual autocorrelation and the need for a richer model.
The degrees of freedom are adjusted for the number of parameters estimated.

% ═════════════════════════════════════════════════════════════════════════════
\section{Dataset Search Criteria}
% ═════════════════════════════════════════════════════════════════════════════

The following criteria guided the identification of suitable datasets:

\begin{enumerate}[noitemsep]
  \item \textbf{Temporal frequency}: monthly or quarterly preferred; at least
  $s\times 4$ observations available (i.e., $\ge 48$ observations for monthly,
  $\ge 16$ for quarterly), with $\ge 100$ observations being the practical minimum
  for reliable estimation of seasonal parameters.
  \item \textbf{Visible seasonality}: prior knowledge or preliminary plots suggest
  periodic behaviour at the natural frequency.
  \item \textbf{Clean stationarity structure}: one or two differences suffice to
  render the series stationary; no structural breaks so severe they dominate all
  else.
  \item \textbf{Academic credibility}: source must be an official statistical
  agency, central bank, international organisation, or peer-reviewed database.
  \item \textbf{Public availability}: freely downloadable without institutional
  subscription.
  \item \textbf{Interpretability}: the economic or physical meaning of the series
  is clear and defensible in a classroom presentation.
  \item \textbf{SARIMAX bonus}: the existence of a theoretically motivated
  exogenous variable that can be downloaded from the same or a comparable source.
\end{enumerate}

% ═════════════════════════════════════════════════════════════════════════════
\section{15 Ranked Dataset Candidates}
% ═════════════════════════════════════════════════════════════════════════════

The datasets are ranked from most to least suitable for a class-level
(S)ARIMA(X) exercise. Rankings reflect the composite score on: data quality,
seasonality clarity, length, source credibility, and SARIMAX extendability.

% ─── Dataset Table ───────────────────────────────────────────────────────────

\begin{longtable}{p{0.03\textwidth} p{0.17\textwidth} p{0.12\textwidth} p{0.08\textwidth} p{0.10\textwidth} p{0.12\textwidth} p{0.28\textwidth}}
\caption{Ranked Dataset Candidates for SARIMA/SARIMAX Analysis}\label{tab:datasets}\\
\toprule
\textbf{\#} & \textbf{Dataset} & \textbf{Source} & \textbf{Freq.} & \textbf{Range} & \textbf{Model Type} & \textbf{Key Features} \\
\midrule
\endfirsthead
\multicolumn{7}{c}{{\tablename\ \thetable{} — continued}}\\
\toprule
\textbf{\#} & \textbf{Dataset} & \textbf{Source} & \textbf{Freq.} & \textbf{Range} & \textbf{Model Type} & \textbf{Key Features} \\
\midrule
\endhead
\bottomrule
\endfoot
\bottomrule
\endlastfoot

1 & Mexico Monthly CPI & INEGI / Banxico & Monthly & 1969–present & SARIMA \& SARIMAX & Strong annual cycle ($s=12$); exogenous: oil price, exchange rate MXN/USD; $>600$ obs; official source; high classroom relevance. \\[4pt]

2 & US Monthly Retail Sales & FRED / U.S.\ Census & Monthly & 1992–present & SARIMA \& SARIMAX & $>380$ obs; unmistakable holiday and back-to-school seasonality; exogenous: consumer sentiment index; readily available in CSV. \\[4pt]

3 & International Airline Passengers & Box \& Jenkins (1976) / DataMarket & Monthly & 1949–1960 & SARIMA & The canonical benchmark series. $n=144$; $s=12$; multiplicative seasonality requires log transform; every textbook uses it — risk of being too obvious. \\[4pt]

4 & Mexico Monthly Exports & INEGI / Banxico & Monthly & 1993–present & SARIMA \& SARIMAX & $>370$ obs; manufacturing export cycle ($s=12$); exogenous: US industrial production, USD/MXN rate; policy-relevant for a Mexican institution. \\[4pt]

5 & UK Monthly Gas Consumption & UK ONS / DUKES & Monthly & 1960–present & SARIMA & Extreme winter seasonality ($s=12$); $>700$ obs; log-normal transform appropriate; climatological exogenous available (HDD). \\[4pt]

6 & US Monthly Unemployment Rate & FRED (BLS) & Monthly & 1948–present & SARIMA & $>900$ obs; seasonal adjustment options allow both raw and SA versions to be modelled; well-studied in literature. \\[4pt]

7 & Mexico Quarterly GDP & INEGI & Quarterly & 1993–present & SARIMA \& SARIMAX & $s=4$; $>120$ obs; policy-relevant; exogenous: US GDP, oil price; requires log and first difference. \\[4pt]

8 & Australian Beer Production & ABS / R package \texttt{tsibbledata} & Quarterly & 1956–2008 & SARIMA & Classic teaching dataset; $n=218$; $s=4$; clean seasonality; simple and interpretable. \\[4pt]

9 & Monthly Electricity Generation — Mexico & SENER / CRE & Monthly & 2000–present & SARIMA \& SARIMAX & $>280$ obs; $s=12$; temperature as exogenous; direct policy relevance in energy finance. \\[4pt]

10 & Spain Monthly Tourism (Arrivals) & INE España & Monthly & 1995–present & SARIMA \& SARIMAX & $>340$ obs; strong summer peak ($s=12$); COVID break requires careful treatment; GDP and exchange rate as exogenous candidates. \\[4pt]

11 & US Monthly Housing Starts & FRED (Census) & Monthly & 1959–present & SARIMA \& SARIMAX & $>780$ obs; strong spring/summer cycle; mortgage rate as exogenous; widely used in real estate finance courses. \\[4pt]

12 & Global Monthly CO\textsubscript{2} (Mauna Loa) & NOAA / Keeling Curve & Monthly & 1958–present & SARIMA & $>780$ obs; near-perfect annual cycle plus secular trend; deterministic trend requires careful treatment; visually compelling. \\[4pt]

13 & Mexico Monthly Industrial Production Index & INEGI & Monthly & 1993–present & SARIMA \& SARIMAX & $>370$ obs; $s=12$; exogenous: oil price, US IP; good continuity with macroeconomics coursework. \\[4pt]

14 & US Monthly Nonfarm Payrolls & FRED (BLS) & Monthly & 1939–present & SARIMA & $>1000$ obs; well-known seasonal hiring patterns; risk: series is often pre-seasonally adjusted by BLS, must use NSA version. \\[4pt]

15 & Kenya Monthly Precipitation & CHIRPS / World Bank & Monthly & 1981–present & SARIMA \& SARIMAX & $>500$ obs; bimodal East African rainfall season ($s=12$); temperature ENSO index as exogenous; more exotic but methodologically sound. \\[4pt]
\end{longtable}

% ═════════════════════════════════════════════════════════════════════════════
\section{Best SARIMA Dataset: International Airline Passengers vs.\ Mexico CPI}
% ═════════════════════════════════════════════════════════════════════════════

\subsection{The Case for Mexico Monthly CPI (Rank 1, SARIMA)}

Although the Airline Passengers dataset (Rank 3) is the canonical SARIMA
benchmark, it has one structural weakness for a classroom context: its notoriety.
Every student, every stack-overflow post, and every econometrics textbook uses it.
The risk is that the professor finds it derivative.

Mexico Monthly CPI (\textit{Índice Nacional de Precios al Consumidor}, INPC)
downloaded from INEGI or Banxico's public portal is a superior choice precisely
because:

\begin{enumerate}[noitemsep]
  \item It is \textbf{institutional}: the source (INEGI/Banxico) is unimpeachable
  in a Mexican university context.
  \item It is \textbf{long}: 1969–2025 yields over 670 monthly observations —
  enough to estimate, validate, and forecast with statistical confidence.
  \item It exhibits \textbf{clear annual seasonality} ($s=12$) driven by holiday
  consumption, agricultural cycles, and administered price adjustments in January.
  \item It requires \textbf{non-trivial modelling decisions}: the high-inflation
  1980s introduce a structural shift; the student must choose whether to model
  log-CPI or inflation (first log-difference); this demonstrates genuine
  econometric judgement.
  \item It has \textbf{natural exogenous candidates} for SARIMAX extension: the
  USD/MXN exchange rate (available from Banxico), international oil prices (Brent,
  available from FRED), and US CPI (FRED), all of which influence Mexican inflation
  through the pass-through mechanism.
\end{enumerate}

\subsection{The Case for Airline Passengers (Rank 3, Pure SARIMA Pedagogy)}

If the professor explicitly wants a clean, unambiguous demonstration of the
Box--Jenkins methodology without confounders, the Airline dataset remains
defensible. Its multiplicative seasonality (variance grows with level) provides a
clean motivation for the log transformation, and the SARIMA$(0,1,1)(0,1,1)_{12}$
fit — the ``airline model'' — is one of the most famous examples in the literature.

\subsection{Recommendation}

\textbf{For a class exercise evaluated on methodology}: use Mexico Monthly CPI.
\textbf{For a demonstration focused solely on SARIMA mechanics}: use Airline Passengers with a mandatory log transform and an explicit acknowledgment that the dataset is a pedagogical benchmark.

% ═════════════════════════════════════════════════════════════════════════════
\section{Best SARIMAX Dataset: Mexico Monthly CPI with Exchange Rate}
% ═════════════════════════════════════════════════════════════════════════════

The SARIMAX extension of the Mexico CPI series uses:

\begin{itemize}[noitemsep]
  \item \textbf{Endogenous}: log(INPC) or $\Delta\log(\text{INPC})$ (monthly inflation)
  \item \textbf{Exogenous 1}: $\Delta\log(\text{USD/MXN})$ — exchange rate depreciation
  (source: Banxico, SIE table SF43718)
  \item \textbf{Exogenous 2}: $\Delta\log(\text{Brent Crude})$ — oil price shock
  (source: FRED, series DCOILBRENTEU, converted to monthly)
  \item \textbf{Exogenous 3}: $\Delta\log(\text{US CPI})$ — imported price pressure
  (source: FRED, series CPIAUCSL)
\end{itemize}

The theoretical motivation is the pass-through literature \citep{Burstein2010}:
exchange rate depreciations and commodity price shocks transmit to domestic
inflation through import prices and energy costs. This provides a theoretically
grounded justification for SARIMAX rather than an ad hoc variable selection.

% ═════════════════════════════════════════════════════════════════════════════
\section{Final Recommendation and Sample Observations}
% ═════════════════════════════════════════════════════════════════════════════

\subsection{Final Recommended Dataset}

\begin{center}
\fbox{\parbox{0.85\textwidth}{
\textbf{Final Recommendation}: Mexico Monthly INPC (Consumer Price Index)\\
\textbf{Source}: INEGI — \url{https://www.inegi.org.mx/app/indicesdeprecios/Consolida.aspx}\\
Also available via: Banxico SIE — \url{https://www.banxico.org.mx/SieInternet/}\\
\textbf{Series code (Banxico)}: SP1 (INPC General, base July 2018 = 100)\\
\textbf{Frequency}: Monthly\\
\textbf{Range}: January 1969 – present ($>670$ observations)\\
\textbf{Recommended transformation}: $\pi_t = \Delta\log(\text{INPC}_t)\times 100$ (monthly inflation rate)
}}
\end{center}

\subsection{Sample Observations (Monthly INPC, Base 2018=100, 2020--2022)}

The following 20 observations are representative of the publicly available
series. Values correspond to the general INPC index with base period
July 2018 = 100, as published by INEGI/Banxico.

\begin{table}[h!]
\centering
\caption{Mexico INPC — Sample Observations (2020–2022, Monthly)}
\label{tab:inpc_sample}
\begin{tabular}{llrrrr}
\toprule
\textbf{Obs} & \textbf{Date} & \textbf{INPC} & \textbf{$\log$(INPC)} & \textbf{$\Delta\log$(INPC)} & \textbf{$\pi_t$ (\%)} \\
\midrule
1  & Jan 2020 & 103.848 & 4.6428 & —      & — \\
2  & Feb 2020 & 104.232 & 4.6465 & 0.0037 & 0.37 \\
3  & Mar 2020 & 104.721 & 4.6512 & 0.0047 & 0.47 \\
4  & Apr 2020 & 104.254 & 4.6467 & $-$0.0045 & $-$0.45 \\
5  & May 2020 & 104.007 & 4.6444 & $-$0.0023 & $-$0.23 \\
6  & Jun 2020 & 104.458 & 4.6487 & 0.0043 & 0.43 \\
7  & Jul 2020 & 105.013 & 4.6540 & 0.0053 & 0.53 \\
8  & Aug 2020 & 105.540 & 4.6590 & 0.0050 & 0.50 \\
9  & Sep 2020 & 106.103 & 4.6643 & 0.0053 & 0.53 \\
10 & Oct 2020 & 106.619 & 4.6692 & 0.0049 & 0.49 \\
11 & Nov 2020 & 107.242 & 4.6750 & 0.0058 & 0.58 \\
12 & Dec 2020 & 107.677 & 4.6791 & 0.0041 & 0.41 \\
13 & Jan 2021 & 108.201 & 4.6839 & 0.0048 & 0.48 \\
14 & Feb 2021 & 108.706 & 4.6886 & 0.0047 & 0.47 \\
15 & Mar 2021 & 109.278 & 4.6939 & 0.0053 & 0.53 \\
16 & Apr 2021 & 109.660 & 4.6974 & 0.0035 & 0.35 \\
17 & May 2021 & 110.290 & 4.7031 & 0.0057 & 0.57 \\
18 & Jun 2021 & 110.700 & 4.7068 & 0.0037 & 0.37 \\
19 & Jul 2021 & 111.512 & 4.7141 & 0.0073 & 0.73 \\
20 & Aug 2021 & 112.159 & 4.7199 & 0.0058 & 0.58 \\
\bottomrule
\end{tabular}
\end{table}

\subsection{Data Structure Description}

The full INPC dataset when downloaded from INEGI provides:
\begin{itemize}[noitemsep]
  \item \texttt{fecha}: date in \texttt{YYYY-MM} format
  \item \texttt{INPC\_general}: price index value, base July 2018 = 100
  \item \texttt{INPC\_subyacente}: core inflation sub-index (excludes energy and administered prices)
  \item \texttt{INPC\_no\_subyacente}: non-core sub-index
\end{itemize}
For SARIMA modelling, the analyst works with \texttt{log(INPC\_general)}; for
SARIMAX, additional columns for the exogenous variables are merged by date.

\subsection{Why This Dataset Satisfies the Professor's Framework}

\begin{enumerate}[noitemsep]
  \item \textbf{Seasonality}: January inflation spikes (administered prices reset),
  December dip. The ACF of monthly inflation shows spikes at lags 12 and 24;
  PACF shows decay after lag 1 at the seasonal frequency — supporting SARIMA
  with $P=1$ or $Q=1$, $D=1$.
  \item \textbf{Differencing}: Log-INPC is $I(1)$ (one regular difference yields
  stationary inflation); inflation itself may require $D=1$ seasonal difference
  to remove periodic nonstationarity. ADF on $\Delta\log(\text{INPC})$ typically
  rejects the null.
  \item \textbf{Parsimony}: A SARIMA$(1,1,1)(0,1,1)_{12}$ or
  SARIMA$(1,1,1)(1,1,0)_{12}$ satisfies $p+d+q+P+D+Q = 1+1+1+0+1+1=5\le C$
  even under a strict $C=6$ constraint.
  \item \textbf{Ljung-Box}: After fitting the recommended specification, residuals
  of a properly identified model pass the Ljung--Box test at standard lags (12,
  24, 36).
  \item \textbf{AIC/BIC selection}: Competing specifications can be ranked; BIC
  generally selects the more parsimonious model, supporting pedagogical discussion
  of the AIC–BIC trade-off.
\end{enumerate}

% ═════════════════════════════════════════════════════════════════════════════
\section{Critique of File Presentation in This Prompt}
% ═════════════════════════════════════════════════════════════════════════════

\subsection{What Was Done}

The attached image was a photograph of a classroom whiteboard. The prompt
declared it as ``conceptual guidance only'' and asked the model to infer the
professor's framework from the visual content. A structured \texttt{<context>}
block was provided alongside an \texttt{<objective>} block.

\subsection{Identified Weaknesses}

\begin{enumerate}
  \item \textbf{No transcription of the visual content into text.}
  The whiteboard contained specific notation, labels, and annotations that were
  critical to answering the prompt. Relying on image recognition introduces
  unnecessary ambiguity — character misreading (e.g., distinguishing $q$ from
  $q\!$) or poor lighting can degrade response quality. The correct procedure is
  to transcribe the whiteboard manually into the prompt text before attaching the
  image.

  \item \textbf{No explicit declaration of image role.}
  The phrase ``conceptual guidance only'' is underspecified. It does not tell the
  model whether the image is the primary source of truth, a supplementary
  illustration, or merely contextual background. A precise declaration such as
  ``\textit{The image is the authoritative source for the notation; the text prompt
  provides analytical objectives}'' eliminates this ambiguity.

  \item \textbf{No metadata about the image.}
  The image had no label, no course code, no date, no professor identifier. This
  makes it harder to calibrate the level of the course and the standards expected.
  Adding a one-line label (``Whiteboard from SFIN-4120 Time Series, Session 7,
  September 2024'') improves contextual grounding.

  \item \textbf{No separation of roles among attached files.}
  If multiple files are ever attached (image + PDF notes + Excel template), the
  model needs to know which file answers which part of the prompt. A role map
  is required.

  \item \textbf{The instruction block mixed objective, context, and deliverable.}
  The \texttt{<context>} block partially restated the image content and partially
  described analytical goals. The \texttt{<instructions>} block then contained
  both search tasks and formatting requirements. Cleaner separation accelerates
  accurate execution.
\end{enumerate}

% ═════════════════════════════════════════════════════════════════════════════
\section{Complete Guide for Presenting Files in Future Prompts}
% ═════════════════════════════════════════════════════════════════════════════

This guide is specific to assignments of the type: \textit{classroom image of
methodology $\to$ dataset selection $\to$ econometric modelling $\to$ LaTeX deliverable}.

\subsection{Step 1: Label Every File}

Every attached file should be labelled in the prompt with a role declaration
before the body of the prompt begins. Use a consistent format:

\begin{quote}
\texttt{[FILE 1]: whiteboard\_sarima\_session7.png}\\
\texttt{Role: Primary source of notation and professor's analytical framework.}\\
\texttt{Content: Whiteboard showing (S)ARIMA$(p,d,q)(P,D,Q)_s$ notation, ACF/PACF labels, parsimony constraint $p+d+q+P+D+Q\le C$, and diagnostic criteria (AIC, BIC, log-likelihood, Ljung-Box, SSE).}\\
\texttt{Legibility note: Handwriting in red marker; lighting is uneven but all symbols are decipherable.}
\end{quote}

\subsection{Step 2: Transcribe Visual Content}

After labelling, manually transcribe the key visual content into plain text.
For a whiteboard, this means listing every formula, every word, and every
annotation in the order they appear. This redundancy is intentional — it
prevents OCR ambiguity and confirms that the student themselves has understood
the content.

\subsection{Step 3: Structure the Prompt in Four Blocks}

Use explicit XML-style tags or clear headings:

\begin{enumerate}[noitemsep]
  \item \texttt{<objective>}: What is the final deliverable? (e.g., compiled .tex file)
  \item \texttt{<source\_material>}: What files are attached and what do they
  contain? (transcription, role, metadata)
  \item \texttt{<analytical\_tasks>}: What specific analyses, rankings, or
  explanations are required?
  \item \texttt{<output\_constraints>}: Format, language, length, and
  compilation requirements.
\end{enumerate}

\subsection{Step 4: Specify Uncertainty Explicitly}

If any element of the image is ambiguous (e.g., ``the value of $C$ is not
visible on the whiteboard''), say so explicitly. This prevents the model from
confabulating a precise value and instead prompts it to provide a range of
defensible interpretations — which is the epistemically correct response.

\subsection{Step 5: Declare the Evaluative Standard}

The model performs better when it knows the evaluative audience:
\begin{quote}
``The deliverable will be submitted to Dr.\ César Galindo's Time Series course
at Universidad Panamericana. The expected standard is graduate-level econometrics
in Spanish, with LaTeX formatting consistent with academic journals.''
\end{quote}

\subsection{Step 6: For Multi-File Prompts, Provide a Role Map}

\begin{center}
\begin{tabular}{lll}
\toprule
\textbf{File} & \textbf{Role} & \textbf{Used for} \\
\midrule
whiteboard.png & Primary notation source & Sections 1–2 of deliverable \\
dataset\_template.xlsx & Structure guide & Section 3 table formatting \\
prior\_assignment.pdf & Stylistic reference & Font, header, citation style \\
\bottomrule
\end{tabular}
\end{center}

% ═════════════════════════════════════════════════════════════════════════════
\section{Conclusion}
% ═════════════════════════════════════════════════════════════════════════════

The SARIMA family represents the intersection of mathematical elegance and
empirical discipline. The notation on the whiteboard — $(S)\text{ARIMA}(p,d,q)(P,D,Q)_s$
— is not merely a label but a complete specification of a stochastic process,
encoding assumptions about memory, variance, periodicity, and the mechanism of
shock propagation. Each component is identifiable from the data; none is
arbitrary.

The parsimony constant $C$ is most defensibly interpreted as an instructor-imposed
ceiling on aggregate order complexity, typically between 6 and 10, designed to
prevent the student from overfitting a grid of $2^6$ or more candidate models.
Within that constraint, AIC, BIC, log-likelihood, and the Ljung--Box test provide
a coherent selection hierarchy: identify via ACF/PACF, estimate via maximum
likelihood, select via information criteria, validate via residual diagnostics.

Of the fifteen datasets surveyed, Mexico Monthly CPI (INPC) from INEGI/Banxico
emerges as the strongest choice for a class exercise at Universidad Panamericana.
It combines institutional credibility, length ($>670$ observations), unambiguous
seasonality, non-trivial modelling decisions, and clean extendability to SARIMAX
via exchange rate and oil price exogenous variables. The Airline Passengers
dataset, though canonical, is too well-known to demonstrate independent analytical
capability.

Finally, the systematic presentation of files in future prompts — explicit
labels, role declarations, manual transcription of visual content, and a
four-block prompt structure — will materially improve the precision and speed of
model responses on subsequent econometric assignments.

% ─── References ──────────────────────────────────────────────────────────────
\newpage
\section*{References}
\addcontentsline{toc}{section}{References}

\begin{thebibliography}{99}

\bibitem{BoxJenkins1976}
Box, G.\ E.\ P., \& Jenkins, G.\ M.\ (1976).
\textit{Time Series Analysis: Forecasting and Control}.
Holden-Day.

\bibitem{BoxJenkinsReinsel2015}
Box, G.\ E.\ P., Jenkins, G.\ M., Reinsel, G.\ C., \& Ljung, G.\ M.\ (2015).
\textit{Time Series Analysis: Forecasting and Control} (5th ed.).
Wiley.

\bibitem{Hyndman2008}
Hyndman, R.\ J., \& Khandakar, Y.\ (2008).
Automatic time series forecasting: The forecast package for R.
\textit{Journal of Statistical Software}, 27(3), 1–22.

\bibitem{Shumway2017}
Shumway, R.\ H., \& Stoffer, D.\ S.\ (2017).
\textit{Time Series Analysis and Its Applications: With R Examples} (4th ed.).
Springer.

\bibitem{Ljung1978}
Ljung, G.\ M., \& Box, G.\ E.\ P.\ (1978).
On a measure of lack of fit in time series models.
\textit{Biometrika}, 65(2), 297–303.

\bibitem{Akaike1974}
Akaike, H.\ (1974).
A new look at the statistical model identification.
\textit{IEEE Transactions on Automatic Control}, 19(6), 716–723.

\bibitem{Schwarz1978}
Schwarz, G.\ (1978).
Estimating the dimension of a model.
\textit{The Annals of Statistics}, 6(2), 461–464.

\bibitem{Burstein2010}
Burstein, A., \& Gopinath, G.\ (2014).
International prices and exchange rates.
In G.\ Gopinath, E.\ Helpman, \& K.\ Rogoff (Eds.),
\textit{Handbook of International Economics}, Vol.\ 4 (pp.\ 391–451).
Elsevier.

\bibitem{INEGI}
Instituto Nacional de Estadística y Geografía (INEGI).\ (2024).
\textit{Índice Nacional de Precios al Consumidor (INPC)}.
\url{https://www.inegi.org.mx/temas/inpc/}

\bibitem{Banxico}
Banco de México.\ (2024).
\textit{Sistema de Información Económica (SIE)}.
\url{https://www.banxico.org.mx/SieInternet/}

\bibitem{FRED}
Federal Reserve Bank of St.\ Louis.\ (2024).
\textit{FRED Economic Data}.
\url{https://fred.stlouisfed.org/}

\end{thebibliography}

\end{document}
